% =============================================================================
%  Case Study: Watchdog Timer (WDT) on Arduino Nano 33 BLE Sense (nRF52840)
%  Author : Akanksh Batra
%  Year   : 2026
%  Compile: pdfLaTeX (Overleaf compatible)
% =============================================================================

\documentclass[12pt, a4paper]{article}

% ── Packages ─────────────────────────────────────────────────────────────────
\usepackage[top=2.5cm, bottom=2.5cm, left=2.5cm, right=2.5cm]{geometry}
\usepackage{fontenc}
\usepackage{inputenc}
\usepackage{microtype}
\usepackage{parskip}
\usepackage{graphicx}
\usepackage{xcolor}
\usepackage{amsmath}
\usepackage{booktabs}
\usepackage{array}
\usepackage{enumitem}

% Code listings
\usepackage{listings}
\definecolor{codebg}{RGB}{245,245,245}
\definecolor{codeframe}{RGB}{200,200,200}
\definecolor{keywordcolor}{RGB}{0,0,180}
\definecolor{commentcolor}{RGB}{60,128,49}
\definecolor{stringcolor}{RGB}{163,21,21}

\lstdefinestyle{arduino}{
  language        = C++,
  backgroundcolor = \color{codebg},
  frame           = single,
  rulecolor       = \color{codeframe},
  basicstyle      = \ttfamily\small,
  keywordstyle    = \color{keywordcolor}\bfseries,
  commentstyle    = \color{commentcolor}\itshape,
  stringstyle     = \color{stringcolor},
  numberstyle     = \tiny\color{gray},
  numbers         = left,
  numbersep       = 8pt,
  stepnumber      = 1,
  breaklines      = true,
  breakatwhitespace = true,
  captionpos      = b,
  tabsize         = 2,
  showstringspaces= false,
  morekeywords    = {uint32_t, bool, void, pinMode, digitalWrite,
                     delay, Serial, HIGH, LOW, OUTPUT, millis,
                     NRF_WDT, NRF_POWER, WDT_CONFIG_HALT_Pause,
                     WDT_CONFIG_HALT_Pos, WDT_CONFIG_SLEEP_Run,
                     WDT_CONFIG_SLEEP_Pos, WDT_RREN_RR0_Msk,
                     WDT_RR_RR_Reload, POWER_RESETREAS_DOG_Msk},
}
\lstset{style=arduino}

% Headers / Footers
\usepackage{fancyhdr}
\pagestyle{fancy}
\fancyhf{}
\rhead{\small Akanksh Batra}
\lhead{\small WDT Case Study — Arduino Nano 33 BLE Sense}
\cfoot{\thepage}
\renewcommand{\headrulewidth}{0.4pt}

% Hyperlinks (load last)
\usepackage[
  colorlinks = true,
  linkcolor  = blue!70!black,
  urlcolor   = blue!70!black,
  citecolor  = blue!70!black,
  pdftitle   = {WDT Case Study — Arduino Nano 33 BLE Sense},
  pdfauthor  = {Akanksh Batra},
]{hyperref}

% ── Title ────────────────────────────────────────────────────────────────────
\title{%
  \textbf{Watchdog Timer (WDT) Case Study}\\[4pt]
  \large Arduino Nano 33 BLE Sense --- nRF52840 SoC\\[2pt]
  \normalsize Bare-Metal Register Programming in Arduino IDE
}
\author{Akanksh Batra}
\date{2026}

% =============================================================================
\begin{document}
% =============================================================================

\maketitle
\thispagestyle{fancy}

\begin{abstract}
This case study documents the configuration and experimental verification of
the hardware Watchdog Timer (WDT) peripheral on the Nordic Semiconductor
nRF52840 System-on-Chip, as deployed on the Arduino Nano 33 BLE Sense
development board. A complete Arduino sketch is presented that (a) programs
the WDT via bare-metal register access, (b) deliberately simulates a firmware
hang after a fixed number of successful WDT feeds, (c) allows the WDT to
trigger an autonomous system reset, and (d) confirms the watchdog reset on
reboot by inspecting the \texttt{POWER→RESETREAS} register. Observed Serial
Monitor output and LED behaviour are reported, and the practical implications
for robust embedded firmware design are discussed.
\end{abstract}

\tableofcontents
\newpage

% ─────────────────────────────────────────────────────────────────────────────
\section{Objective}
% ─────────────────────────────────────────────────────────────────────────────

The primary objective of this case study is to:

\begin{enumerate}[label=\arabic*.]
  \item Understand the architectural role of a hardware Watchdog Timer in
        embedded systems.
  \item Configure the nRF52840 WDT peripheral using bare-metal C++ register
        writes inside the Arduino IDE environment.
  \item Observe and verify the automatic recovery behaviour when the WDT timeout
        expires due to a simulated firmware lockup.
  \item Demonstrate how to distinguish a WDT-induced reset from other reset
        sources using the \texttt{POWER→RESETREAS} register.
\end{enumerate}

% ─────────────────────────────────────────────────────────────────────────────
\section{Introduction}
% ─────────────────────────────────────────────────────────────────────────────

Embedded systems deployed in unattended or safety-critical environments must
be capable of recovering autonomously from software faults such as infinite
loops, deadlocks, or stack overflows. A \textbf{Watchdog Timer (WDT)} is a
hardware peripheral that fulfils this requirement: it maintains an independent
countdown timer that triggers a system reset unless the application
periodically reloads (``feeds'' or ``pets'') it within a programmed deadline.

The Nordic Semiconductor nRF52840, the SoC at the heart of the Arduino Nano 33
BLE Sense, features a dedicated WDT peripheral clocked by the 32.768\,kHz
Low-Frequency Clock (LFCLK). This clock source is independent of the CPU and
peripheral bus clocks, ensuring that the WDT continues to operate even when the
main system is locked up. Once started, the nRF52840 WDT \emph{cannot be
stopped without a device reset}, providing a strong safety guarantee.

% ─────────────────────────────────────────────────────────────────────────────
\section{Hardware and Software Requirements}
% ─────────────────────────────────────────────────────────────────────────────

\subsection{Hardware}

\begin{table}[h!]
  \centering
  \begin{tabular}{>{\bfseries}p{4cm} p{9cm}}
    \toprule
    Component       & Details \\
    \midrule
    Development Board & Arduino Nano 33 BLE Sense (Rev1 or Rev2) \\
    SoC               & Nordic Semiconductor nRF52840
                        (64\,MHz Cortex-M4F, 1\,MB Flash, 256\,kB RAM) \\
    Connection        & USB-A to Micro-USB cable (power + Serial CDC) \\
    \bottomrule
  \end{tabular}
  \caption{Hardware components used in this case study.}
\end{table}

\subsection{Software}

\begin{table}[h!]
  \centering
  \begin{tabular}{>{\bfseries}p{4cm} p{9cm}}
    \toprule
    Tool              & Version / Details \\
    \midrule
    Arduino IDE       & 2.x (tested on 2.3+) \\
    Board Package     & Arduino Mbed OS Nano Boards $\geq$ 4.0.0 \\
    Header            & \texttt{<nrf.h>} (bundled with core) \\
    Serial Monitor    & 115\,200 baud, Newline line-ending \\
    \bottomrule
  \end{tabular}
  \caption{Software environment for this case study.}
\end{table}

No external libraries are required. Direct register access is provided by the
\texttt{<nrf.h>} CMSIS header shipped with the Mbed OS core.

% ─────────────────────────────────────────────────────────────────────────────
\section{Working Principle of the nRF52840 WDT}
% ─────────────────────────────────────────────────────────────────────────────

\subsection{Clock Source and Timeout Formula}

The WDT peripheral is driven by the 32.768\,kHz LFCLK oscillator. The
\textbf{Counter Reload Value (CRV)} register programs the timeout period
according to:

\[
  \text{CRV} = \frac{T_{\mathrm{timeout}}[\mathrm{ms}] \times 32{,}768}{1000}
\]

For a 3-second timeout:
\[
  \text{CRV} = \frac{3000 \times 32{,}768}{1000} = 98{,}304
\]

\subsection{Key Registers}

\begin{table}[h!]
  \centering
  \begin{tabular}{>{\ttfamily}l p{9cm}}
    \toprule
    \textrm{Register}           & Description \\
    \midrule
    NRF\_WDT->CONFIG            & Controls WDT behaviour during CPU halt
                                  (debug) and sleep. \\
    NRF\_WDT->CRV               & Counter Reload Value — sets the timeout
                                  period. \\
    NRF\_WDT->RREN              & Reload Register Enable — selects which
                                  RR[n] registers must be written to pet
                                  the WDT. \\
    NRF\_WDT->TASKS\_START      & Write 1 to start the WDT (irreversible
                                  without reset). \\
    NRF\_WDT->RR[0]             & Reload Register 0 — write
                                  \texttt{WDT\_RR\_RR\_Reload} (0x6E524635)
                                  to reset the counter. \\
    NRF\_POWER->RESETREAS       & Bitmask of last reset cause; bit
                                  \texttt{DOG} indicates a WDT reset. \\
    \bottomrule
  \end{tabular}
  \caption{nRF52840 WDT and POWER registers used in this case study.}
\end{table}

\subsection{Operational Flow}

\begin{enumerate}[label=\arabic*.]
  \item \textbf{Configure} --- Write \texttt{CONFIG}, \texttt{CRV}, and
        \texttt{RREN} before starting.
  \item \textbf{Start} --- Write 1 to \texttt{TASKS\_START}. The WDT
        counter begins counting up.
  \item \textbf{Feed} --- Application writes \texttt{WDT\_RR\_RR\_Reload} to
        \texttt{RR[0]} before the counter reaches \texttt{CRV}, resetting it
        to zero.
  \item \textbf{Timeout} --- If the counter reaches \texttt{CRV} without a
        feed, the WDT asserts a system reset.
  \item \textbf{Detect} --- On reboot, firmware checks \texttt{RESETREAS.DOG}
        to determine if a WDT reset occurred, then clears the register.
\end{enumerate}

% ─────────────────────────────────────────────────────────────────────────────
\section{Code Listing}
% ─────────────────────────────────────────────────────────────────────────────

The complete Arduino sketch is reproduced below with inline commentary.

\subsection{WDT Configuration and Feed Functions}

\begin{lstlisting}[caption={WDT configuration (\texttt{configureWDT}) and feed
  (\texttt{feedWDT}) functions.}]
#include <nrf.h>

#define WDT_TIMEOUT_MS   3000
#define FEED_INTERVAL_MS 1000
#define HANG_AFTER_FEEDS 10

uint32_t feedCount = 0;

void configureWDT(uint32_t timeoutMs) {
  // Pause WDT on debug halt; keep running during CPU sleep
  NRF_WDT->CONFIG = (WDT_CONFIG_HALT_Pause << WDT_CONFIG_HALT_Pos)
                  | (WDT_CONFIG_SLEEP_Run  << WDT_CONFIG_SLEEP_Pos);
  // Set timeout: CRV = (ms * 32768) / 1000
  NRF_WDT->CRV = (timeoutMs * 32768ULL) / 1000;
  // Enable reload register 0
  NRF_WDT->RREN |= WDT_RREN_RR0_Msk;
  // Start WDT (cannot be stopped without reset)
  NRF_WDT->TASKS_START = 1;
}

void feedWDT() {
  NRF_WDT->RR[0] = WDT_RR_RR_Reload;
}
\end{lstlisting}

\subsection{Reset Reason Detection}

\begin{lstlisting}[caption={Reading and clearing the \texttt{RESETREAS}
  register.}]
bool wasResetByWDT() {
  return (NRF_POWER->RESETREAS & POWER_RESETREAS_DOG_Msk) != 0;
}

void clearResetReason() {
  // Write-1-to-clear: write current value back to clear all set bits
  NRF_POWER->RESETREAS = NRF_POWER->RESETREAS;
}
\end{lstlisting}

\subsection{setup() and loop()}

\begin{lstlisting}[caption={\texttt{setup()} initialises the WDT and reports
  the boot reason; \texttt{loop()} feeds the WDT and simulates a hang.}]
void setup() {
  Serial.begin(115200);
  while (!Serial && millis() < 3000);
  pinMode(LED_BUILTIN, OUTPUT);

  Serial.println("=========================================");
  Serial.println(" Nano 33 BLE Sense - WDT Case Study Demo ");
  Serial.println("=========================================");

  if (wasResetByWDT()) {
    Serial.println(">>> SYSTEM RECOVERED FROM A WATCHDOG RESET! <<<");
    Serial.println(">>> WDT functionality CONFIRMED working.    <<<");
    for (int i = 0; i < 10; i++) {
      digitalWrite(LED_BUILTIN, HIGH); delay(100);
      digitalWrite(LED_BUILTIN, LOW);  delay(100);
    }
  } else {
    Serial.println("Normal power-on / manual reset detected.");
  }

  clearResetReason();
  configureWDT(WDT_TIMEOUT_MS);
  Serial.println("WDT configured with 3000 ms timeout.");
  Serial.println("Feeding WDT every 1000 ms...");
}

void loop() {
  feedWDT();
  feedCount++;

  digitalWrite(LED_BUILTIN, HIGH);
  delay(150);
  digitalWrite(LED_BUILTIN, LOW);

  Serial.print("WDT fed. Count = ");
  Serial.println(feedCount);

  if (feedCount >= HANG_AFTER_FEEDS) {
    Serial.println("Simulating firmware hang (infinite loop)...");
    Serial.println("WDT will NOT be fed from this point onward.");
    Serial.println("System will auto-reset in ~3 seconds.");
    Serial.flush();
    while (1) { /* deliberate hang — no feedWDT() */ }
  }

  delay(FEED_INTERVAL_MS - 150);
}
\end{lstlisting}

% ─────────────────────────────────────────────────────────────────────────────
\section{Expected Serial Monitor Output}
% ─────────────────────────────────────────────────────────────────────────────

\begin{lstlisting}[language={}, caption={Serial Monitor output showing the
  full hang-reset-recovery cycle.}, basicstyle=\ttfamily\footnotesize]
=========================================
 Nano 33 BLE Sense - WDT Case Study Demo
=========================================
Normal power-on / manual reset detected.
WDT configured with 3000 ms timeout.
Feeding WDT every 1000 ms...
-----------------------------------------
WDT fed. Count = 1
WDT fed. Count = 2
WDT fed. Count = 3
WDT fed. Count = 4
WDT fed. Count = 5
WDT fed. Count = 6
WDT fed. Count = 7
WDT fed. Count = 8
WDT fed. Count = 9
WDT fed. Count = 10
-----------------------------------------
Simulating firmware hang (infinite loop)...
WDT will NOT be fed from this point onward.
System will auto-reset in ~3 seconds.

--- [ ~3 s pause ... WDT fires ... board resets ] ---

=========================================
 Nano 33 BLE Sense - WDT Case Study Demo
=========================================
>>> SYSTEM RECOVERED FROM A WATCHDOG RESET! <<<
>>> WDT functionality CONFIRMED working.    <<<
WDT configured with 3000 ms timeout.
Feeding WDT every 1000 ms...
-----------------------------------------
WDT fed. Count = 1
...
\end{lstlisting}

% ─────────────────────────────────────────────────────────────────────────────
\section{Observations}
% ─────────────────────────────────────────────────────────────────────────────

\begin{enumerate}[label=\arabic*.]
  \item \textbf{Normal Operation (feeds 1–10):}  
        The LED produces a clean single blink per second. Each feed resets the
        WDT counter well within the 3-second window. No reset occurs.

  \item \textbf{Simulated Hang:}  
        After feed \#10, the sketch enters \texttt{while(1)} and the LED
        freezes in the OFF state. The Serial Monitor shows the last three
        diagnostic messages then falls silent.

  \item \textbf{WDT-Triggered Reset ($\approx$3\,s after hang):}  
        The board resets autonomously without any user interaction. The USB-CDC
        Serial connection briefly drops and reconnects.

  \item \textbf{Recovery Detection:}  
        On the second boot, \texttt{POWER→RESETREAS.DOG} is set. The sketch
        correctly identifies the WDT reset and emits the recovery banner.

  \item \textbf{Rapid Blink:}  
        The LED blinks 10 times at 100\,ms intervals — a distinctive pattern
        that visually confirms autonomous fault recovery.

  \item \textbf{Cycle Repeats:}  
        Because the sketch always simulates a hang after 10 feeds, the
        hang-reset-recovery cycle continues indefinitely, making this an
        effective live demonstration.
\end{enumerate}

% ─────────────────────────────────────────────────────────────────────────────
\section{Statement of Results}
% ─────────────────────────────────────────────────────────────────────────────

The experiment conclusively demonstrates that:

\begin{itemize}
  \item The nRF52840 hardware WDT can be configured and exercised through
        bare-metal register writes inside the Arduino Mbed OS environment
        without any third-party libraries.
  \item A 3-second WDT timeout fires reliably when the application stops
        feeding the WDT, producing an autonomous system reset.
  \item The \texttt{POWER→RESETREAS} register accurately reports the WDT
        reset cause, enabling firmware to implement graceful post-fault
        recovery logic.
\end{itemize}

% ─────────────────────────────────────────────────────────────────────────────
\section{Real-World Applications}
% ─────────────────────────────────────────────────────────────────────────────

\begin{description}
  \item[Remote IoT Nodes] Sensors deployed in remote locations cannot be
    manually reset. A WDT ensures they recover from transient software faults
    and resume data transmission without human intervention.

  \item[Medical Devices] Regulatory standards (IEC 62304) mandate that
    safety-critical software include fault-detection mechanisms. A hardware WDT
    provides a last-resort recovery independent of the software stack.

  \item[Industrial Automation] PLCs and motion controllers must continue
    operating (or fail safely) despite software errors. WDTs guard against
    lockups in time-critical control loops.

  \item[Wearables and Consumer Electronics] Battery-powered devices that
    experience rare software hangs can recover silently, improving perceived
    reliability and user experience.

  \item[Aerospace and Automotive (ASIL/DAL-rated firmware)] WDTs are a
    mandatory safety mechanism in automotive (ISO 26262) and aerospace systems
    to meet fault-tolerance requirements.
\end{description}

% ─────────────────────────────────────────────────────────────────────────────
\section{Conclusion}
% ─────────────────────────────────────────────────────────────────────────────

This case study successfully demonstrated the end-to-end operation of the
hardware Watchdog Timer on the nRF52840 SoC using the Arduino Nano 33 BLE
Sense. Through direct register manipulation — without relying on high-level
Arduino abstractions — the WDT was configured with a 3-second timeout,
periodically fed during normal operation, and proven to trigger an autonomous
reset when the simulated firmware hang prevented further feeding. Post-reset
detection via \texttt{POWER→RESETREAS} confirmed the exact cause of the reset,
illustrating a complete fault-detection and recovery cycle.

The technique presented here is directly applicable to production embedded
firmware. A 3-line addition to any \texttt{loop()} and a startup check of
\texttt{RESETREAS} can dramatically improve the resilience of any nRF52840-based
product with minimal code overhead.

\vspace{1em}
\noindent\rule{\linewidth}{0.4pt}

\begin{center}
  \textit{Case study prepared by \textbf{Akanksh Batra} --- 2026}\\
  Repository: \url{https://github.com/ak4204/nano33ble-wdt-case-study}
\end{center}

\end{document}
